\documentclass[11pt]{article}
\usepackage{fontspec,amsmath,amssymb,amsthm,hyperref}
\usepackage[margin=1in]{geometry}
\setlength{\emergencystretch}{3em}
\title{Bounded radial weights approximate every square-integrable initial vector}
\author{LangevinLean mathematical exposition}
\date{}
\begin{document}\maketitle
\noindent Source module: \texttt{AverageLMC/LaplacianWeightClipping.lean}. Compile this proof document with XeLaTeX or LuaLaTeX. The accompanying HTML pairs the mathematical argument with the exact Lean statements and records the scope of the formalization.


The initial coupling error in the Laplacian proof is only square integrable. This module constructs bounded radial clippings whose L2 norms never exceed the original norm. Their pairings converge against every fixed L2 vector by dominated convergence, allowing the later compact-test argument to use bounded scalar weights without adding boundedness to the main theorem hypotheses.



The measurable radial clipping has \(c_n(0)=0\); for \(x\ne0\) it uses radius \(n+1\):


\[c_n(x)=\min\!\left\{1,\frac{n+1}{\|x\|}\right\}x,
\qquad\|c_n(x)\|\le\min\{\|x\|,n+1\}.\]


Clipping never increases the actual \(L^2\) norm and converges against every square-integrable vector:


\[\|c_n(\Psi)\|_{L^2(P)}\le\|\Psi\|_{L^2(P)},\qquad
\langle c_n(\Psi),D\rangle_{L^2(P)}
\longrightarrow\langle\Psi,D\rangle_{L^2(P)}.\]
\section{Define radial clipping and prove both norm bounds}


Clip a vector x by multiplying it by min(1,(n+1)/norm(x)). The coefficient lies between zero and one, so the clipped norm is at most norm(x). It is also at most n+1: at a nonzero vector cancel the positive norm against the radial quotient, and handle the zero vector directly. These two bounds serve different purposes: the radius makes every coordinate a bounded test weight, while domination by the original vector supplies the integrable majorant and preserves the sharp final L2 norm.



\noindent Lean declarations: \texttt{clipWeight}, \texttt{norm\_clipWeight\_le}, \texttt{norm\_clipWeight\_le\_radius}.

\section{Obtain eventual equality and measurability}


If norm(x) is at most n+1, the minimum coefficient is one and clipping fixes x. For a fixed vector choose a natural number above its norm. Every larger index then fixes that vector exactly, giving eventual pointwise equality with no approximation-rate premise. Norm, real division, minimum and scalar multiplication are measurable, so the clipping map is measurable even with the real quotient convention at zero.



\noindent Lean declarations: \texttt{clipWeight\_eq\_of\_norm\_le}, \texttt{clipWeight\_eventually\_eq}, \texttt{measurable\_clipWeight}.

\section{Lift the clippings into the genuine L2 space}


Compose the measurable clipping with an L2 representative. The pointwise norm domination by the original representative proves MemLp at exponent two. The resulting canonical toLp element is lpClippedWeight, with its almost-everywhere representative identity proved explicitly. Monotonicity of the L2 norm under the same pointwise domination gives norm(lpClippedWeight X n) at most norm(X), uniformly in n.



\noindent Lean declarations: \texttt{clippedWeight\_memLp}, \texttt{lpClippedWeight}, \texttt{lpClippedWeight\_coe}, \texttt{lpClippedWeight\_norm\_le}.

\section{Pass every fixed L2 pairing through the clipping limit}


For X and D in L2, Holder gives integrability of norm(X) times norm(D). Cauchy-Schwarz bounds the absolute pointwise inner product of the clipping of X with D by this common majorant. Eventual pointwise equality gives pointwise convergence of the integrands, and dominated convergence gives convergence of their actual integrals. Rewriting the Hilbert-space L2 inner product as that integral and using the almost-everywhere clipping representative proves convergence of the exact L2 pairings.



\noindent Lean declarations: \texttt{lpClippedWeight\_inner\_tendsto}.

\end{document}
